%=====================================================================
\section{Introduction}
\label{sec:intro}

\subsection{The optimisation problem in one page}

During a planning period $[0,T]$ a set $I=\{1,\dots,n\}$ of parked electric vehicles (EVs) raise charging requests. Request $i$ is released at minute $r_i$ and needs $E_i$ kWh. It must be served by exactly one resource:
\begin{itemize}
  \item a \textbf{gun} $g$ of a fixed charging station (FCS) $f(g)$ within the EV's remaining range $R_i$: the EV drives there ($d_{if}$ km, $t_{if}$ minutes), queues for at most $W_i$ minutes and charges for $\tau^\Fcs_i$ minutes; or
  \item a \textbf{mobile charging van} (MCS) $m$: the van comes from the depot or from its previous EV, the EV waits $q^\Mcs_i$ minutes and is charged for $\tau^\Mcs_i<\tau^\Fcs_i$ minutes while the driver is free.
\end{itemize}
A van serves at most $Q$ EVs, carries $B$ kWh for both driving ($\beta$ kWh/km) and delivering ($E_i/\eta$), and must be back at the depot by $T$. Every EV must be fully charged by $T$.

Every EV carries an income index $\theta_i\in[0,1]$. It sets the customer's money sensitivity $\mu_i=\mu(\theta_i)$ (non-increasing), values of time $c^{a}_i=c^{a}(\theta_i)$ for driving, queueing, charging and van service $a\in\{\mathrm D,\mathrm Q,\mathrm C,\mathrm V\}$ (non-decreasing), and queue tolerance $W_i=W(\theta_i)$ (non-increasing). The objective of formulation 3 adds the operator's cost and the cost customers perceive:
\begin{equation}
Z=\underbrace{Z_1}_{\text{EV driving}}+\underbrace{Z_2}_{\text{van km}}+\underbrace{Z_3}_{\text{dispatch}}+\underbrace{Z_4}_{\text{energy}}+\underbrace{Z_5}_{\text{tariff}}+\underbrace{Z_6}_{\text{time}} .
\label{eq:Z}
\end{equation}
The exact terms are in \code{main\_formulation\_3.pdf}, eqs.\ (1)--(23); Section~\ref{sec:mdp} restates the parts the RL environment needs. The decisions are: who serves each EV, in what order each gun and each van takes its EVs, and when each session starts.

\subsection{Why reinforcement learning at all}

The exact model is a mixed-integer linear program (MILP). The project solves it with Gurobi on 560 EV-Solomon instances (56 Solomon files, $n=10,20,\dots,100$) with a 300-second limit per instance. That run gives a reference of high quality, but it also shows the limits of exact solving:
\begin{itemize}
  \item mean solve time 150.7\,s per instance, and only 256 of 560 instances proven optimal;
  \item 30 instances end the time limit \emph{without any plan} (R1/RC1 with $n\ge70$ and one C2 instance), because the gun and queue constraints make even a first feasible solution hard to find;
  \item every new instance starts from scratch: nothing is learned from the 559 others.
\end{itemize}
A learned \emph{constructive heuristic} answers exactly these points. A neural network is trained once, offline, on many random instances drawn from the same rules as the dataset. Afterwards it builds a full plan for a new instance in a fraction of a second, it can be asked for many alternative plans in parallel, and it is not stopped by the difficulty of proving optimality. The price is that it gives no optimality guarantee: its plans are judged by comparing them with the MILP.

\subsection{What was built}

\begin{description}
  \item[RL1] (\code{RL1\_reinforcement\_learning\_method\_1/}) an attention-based policy that builds the plan one decision at a time. Each decision is a pair \emph{(EV $i$, resource $k$)}. A simulator of formulation 3 computes start times, checks feasibility, and returns the exact increase of $Z$ as the (negative) reward. Training is REINFORCE with the POMO shared baseline after a short imitation warm-up. At test time 520 plans per instance are built (8 map symmetries $\times$ 65 rollouts) and the best is kept.
  \item[RL2] (\code{RL2\_reinforcement\_learning\_method\_2/}) the same method with three changes in training (bigger batches, more large instances, a larger validation set) and \emph{Efficient Active Search}: for each test instance the trained network is frozen and the instance's embeddings are tuned on that instance for up to 10 seconds. RL1 is not changed by RL2; both read the same dataset.
\end{description}
Both write their plans in the project's common results format, and every plan is scored by the same checker (\code{check\_plan\_f3}) as the MILP and greedy plans, so the comparison is exact.

\subsection{How to read this report}

Section~\ref{sec:glossary} defines the RL vocabulary. Section~\ref{sec:chain} tells how the method was chosen, step by step, including the alternatives that were rejected. Section~\ref{sec:papers} lists the papers and what was taken from each. Section~\ref{sec:assumptions} collects every assumption with its reason and its effect. Sections~\ref{sec:mdp}--\ref{sec:inference} give the method itself: environment, network, training, and test-time search. Section~\ref{sec:principles} states the principles the implementation follows. Section~\ref{sec:evidence} gives the evidence so far, Section~\ref{sec:limits} the limits and next steps, and Section~\ref{sec:repro} how to reproduce everything.
